\documentclass{amsart}
% For dealing with references we use the comment environment
\usepackage{verbatim}
\newenvironment{reference}{\comment}{\endcomment}
%\newenvironment{reference}{}{}
\newenvironment{slogan}{\comment}{\endcomment}
\newenvironment{history}{\comment}{\endcomment}

% For commutative diagrams we use Xy-pic
\usepackage[all]{xy}

% We use 2cell for 2-commutative diagrams.
\xyoption{2cell}
\UseAllTwocells

% We use multicol for the list of chapters between chapters
\usepackage{multicol}

% This is generally recommended for better output
\usepackage{lmodern}
\usepackage[T1]{fontenc}

% For cross-file-references

% Package for hypertext links:
\usepackage{hyperref}

% For any local file, say "hello.tex" you want to link to please

% Theorem environments.
%
\theoremstyle{plain}
\newtheorem{theorem}[subsection]{Theorem}
\newtheorem{proposition}[subsection]{Proposition}
\newtheorem{lemma}[subsection]{Lemma}

\theoremstyle{definition}
\newtheorem{definition}[subsection]{Definition}
\newtheorem{example}[subsection]{Example}
\newtheorem{exercise}[subsection]{Exercise}
\newtheorem{situation}[subsection]{Situation}

\theoremstyle{remark}
\newtheorem{remark}[subsection]{Remark}
\newtheorem{remarks}[subsection]{Remarks}

\numberwithin{equation}{subsection}

% Macros
%
\def\lim{\mathop{\mathrm{lim}}\nolimits}
\def\colim{\mathop{\mathrm{colim}}\nolimits}
\def\Spec{\mathop{\mathrm{Spec}}}
\def\Hom{\mathop{\mathrm{Hom}}\nolimits}
\def\Ext{\mathop{\mathrm{Ext}}\nolimits}
\def\SheafHom{\mathop{\mathcal{H}\!\mathit{om}}\nolimits}
\def\SheafExt{\mathop{\mathcal{E}\!\mathit{xt}}\nolimits}
\def\Sch{\mathit{Sch}}
\def\Mor{\mathop{\mathrm{Mor}}\nolimits}
\def\Ob{\mathop{\mathrm{Ob}}\nolimits}
\def\Sh{\mathop{\mathit{Sh}}\nolimits}
\def\NL{\mathop{N\!L}\nolimits}
\def\CH{\mathop{\mathrm{CH}}\nolimits}
\def\proetale{{pro\text{-}\acute{e}tale}}
\def\etale{{\acute{e}tale}}
\def\QCoh{\mathit{QCoh}}
\def\Ker{\mathop{\mathrm{Ker}}}
\def\Im{\mathop{\mathrm{Im}}}
\def\Coker{\mathop{\mathrm{Coker}}}
\def\Coim{\mathop{\mathrm{Coim}}}

% Boxtimes
%
\DeclareMathSymbol{\boxtimes}{\mathbin}{AMSa}{"02}

%
% Macros for moduli stacks/spaces
%
\def\QCohstack{\mathcal{QC}\!\mathit{oh}}
\def\Cohstack{\mathcal{C}\!\mathit{oh}}
\def\Spacesstack{\mathcal{S}\!\mathit{paces}}
\def\Quotfunctor{\mathrm{Quot}}
\def\Hilbfunctor{\mathrm{Hilb}}
\def\Curvesstack{\mathcal{C}\!\mathit{urves}}
\def\Polarizedstack{\mathcal{P}\!\mathit{olarized}}
\def\Complexesstack{\mathcal{C}\!\mathit{omplexes}}
% \Pic is the operator that assigns to X its picard group, usage \Pic(X)
% \Picardstack_{X/B} denotes the Picard stack of X over B
% \Picardfunctor_{X/B} denotes the Picard functor of X over B
\def\Pic{\mathop{\mathrm{Pic}}\nolimits}
\def\Picardstack{\mathcal{P}\!\mathit{ic}}
\def\Picardfunctor{\mathrm{Pic}}
\def\Deformationcategory{\mathcal{D}\!\mathit{ef}}

\setlength{\emergencystretch}{3em}
\begin{document}
\title[Stacks Project, Tag 00UE]{1325 --- Étale ring maps over arbitrary bases are standard étale after a principal localization at every étale prime}
\maketitle
\noindent Copyright (C) 2005--2025 Johan de Jong. Stacks Project authors. Source: \href{https://stacks.math.columbia.edu/tag/00UE}{Tag 00UE}.

\noindent Editorial difficulty note (not author text). Difficulty H2: The proof isolates a separable residue-field factor, constructs a monogenic finite model via a primitive residue generator, and manufactures a monic polynomial with an invertible derivative by combining two kernel polynomials. The full eleven-stage finite-localization and polynomial construction is substantially beyond qualification étale definitions..

\noindent Editorial provenance note (not author text). History: excerpt from revision \texttt{a0444\allowbreak 6e57e\allowbreak c1fbc\allowbreak 252a8\allowbreak 71afc\allowbreak ec775\allowbreak 2fb28\allowbreak 07b14}, algebra.tex, lines 39843--40062. Reference presentation was adapted; original-author context and core support blocks were retained or added. The mathematical wording is unchanged. Licensed under GNU FDL 1.2 or later; no invariant sections or cover texts. See ../licenses/COPYING. Context blocks reproduce author definitions and supporting results; they are not additional counted problems.

\begin{definition}
\label{definition-etale}
Let $R \to S$ be a ring map. We say $R \to S$ is {\it \'etale} if it is
of finite presentation and the naive cotangent complex
$\NL_{S/R}$ is quasi-isomorphic to zero: this means that $H_1(\NL_{S/R}) = 0$
and $\Omega_{S/R} = 0$. Given a prime $\mathfrak q$
of $S$ we say that $R \to S$ is {\it \'etale at $\mathfrak q$}
if there exists a $g \in S$, $g \not \in \mathfrak q$ such that
$R \to S_g$ is \'etale.
\end{definition}

\begin{definition}
\label{definition-standard-etale}
Let $R$ be a ring. Let $g , f  \in R[x]$.
Assume that $f$ is monic and the derivative $f'$ is invertible in
the localization $R[x]_g/(f)$.
In this case the ring map $R \to R[x]_g/(f)$ is said to be
{\it standard \'etale}.
\end{definition}

\begin{lemma}
\label{lemma-etale}
Results on \'etale ring maps.
\begin{enumerate}
\item The ring map $R \to R_f$ is \'etale for any ring $R$ and any $f \in R$.
\item Compositions of \'etale ring maps are \'etale.
\item A base change of an \'etale ring map is \'etale.
\item The property of being \'etale is local: Given a ring map
$R \to S$ and elements $g_1, \ldots, g_m \in S$ which generate the unit ideal
such that $R \to S_{g_j}$ is \'etale for $j = 1, \ldots, m$ then
$R \to S$ is \'etale.
\item Given $R \to S$ of finite presentation, and a flat ring map
$R \to R'$, set $S' = R' \otimes_R S$. The set of primes where $R' \to S'$
is \'etale is the inverse image via $\Spec(S') \to \Spec(S)$
of the set of primes where $R \to S$ is \'etale.
\item An \'etale ring map is syntomic, in particular flat.
\item If $S$ is finite type over a field $k$, then $S$ is \'etale over
$k$ if and only if $\Omega_{S/k} = 0$.
\item Any \'etale ring map $R \to S$ is the base change of an \'etale
ring map $R_0 \to S_0$ with $R_0$ of finite type over $\mathbf{Z}$.
\item Let $A = \colim A_i$ be a filtered colimit of rings.
Let $A \to B$ be an \'etale ring map. Then there exists an \'etale ring
map $A_i \to B_i$ for some $i$ such that $B \cong A \otimes_{A_i} B_i$.
\item Let $A$ be a ring. Let $S$ be a multiplicative subset of $A$.
Let $S^{-1}A \to B'$ be \'etale. Then there exists an \'etale ring map
$A \to B$ such that $B' \cong S^{-1}B$.
\item Let $A$ be a ring. Let $B = B' \times B''$ be a product of $A$-algebras.
Then $B$ is \'etale over $A$ if and only if both $B'$ and $B''$ are
\'etale over $A$.
\end{enumerate}
\end{lemma}

\begin{proof}
In each case we use the corresponding result for smooth ring maps with
a small argument added to show that $\Omega_{S/R}$ is zero.

\medskip\noindent
Proof of (1). The ring map $R \to R_f$ is smooth and $\Omega_{R_f/R} = 0$.

\medskip\noindent
Proof of (2). The composition $A \to C$ of smooth maps $A \to B$ and
$B \to C$ is smooth, see Lemma \href{https://stacks.math.columbia.edu/tag/00TD}{lemma compose smooth (Tag 00TD)}. By
Lemma \href{https://stacks.math.columbia.edu/tag/00RS}{lemma exact sequence differentials (Tag 00RS)} we see that
$\Omega_{C/A}$ is zero as both $\Omega_{C/B}$ and $\Omega_{B/A}$ are zero.

\medskip\noindent
Proof of (3). Let $R \to S$ be \'etale and $R \to R'$ be arbitrary.
Then $R' \to S' = R' \otimes_R S$ is smooth, see
Lemma \href{https://stacks.math.columbia.edu/tag/00T4}{lemma base change smooth (Tag 00T4)}. Since
$\Omega_{S'/R'} = S' \otimes_S \Omega_{S/R}$ by
Lemma \href{https://stacks.math.columbia.edu/tag/00RV}{lemma differentials base change (Tag 00RV)}
we conclude that $\Omega_{S'/R'} = 0$. Hence $R' \to S'$ is \'etale.

\medskip\noindent
Proof of (4). Assume the hypotheses of (4). By
Lemma \href{https://stacks.math.columbia.edu/tag/00TC}{lemma locally smooth (Tag 00TC)} we see that $R \to S$ is smooth.
We are also given that $\Omega_{S_{g_i}/R} = (\Omega_{S/R})_{g_i} = 0$
for all $i$. Then $\Omega_{S/R} = 0$, see Lemma \href{https://stacks.math.columbia.edu/tag/00EO}{lemma cover (Tag 00EO)}.

\medskip\noindent
Proof of (5). The result for smooth maps is
Lemma \href{https://stacks.math.columbia.edu/tag/00TG}{lemma flat base change locus smooth (Tag 00TG)}.
In the proof of that lemma we used that $\NL_{S/R} \otimes_S S'$
is homotopy equivalent to $\NL_{S'/R'}$.
This reduces us to showing that if $M$ is a finitely presented
$S$-module the set of primes $\mathfrak q'$ of $S'$
such that $(M \otimes_S S')_{\mathfrak q'} = 0$ is the inverse
image of the set of primes $\mathfrak q$ of $S$ such that
$M_{\mathfrak q} = 0$. This follows from Lemma \href{https://stacks.math.columbia.edu/tag/0BUR}{lemma support base change (Tag 0BUR)}.

\medskip\noindent
Proof of (6). Follows directly from the corresponding result for
smooth ring maps (Lemma \href{https://stacks.math.columbia.edu/tag/00TA}{lemma smooth syntomic (Tag 00TA)}).

\medskip\noindent
Proof of (7). Follows from Lemma \href{https://stacks.math.columbia.edu/tag/00TT}{lemma characterize smooth over field (Tag 00TT)}
and the definitions.

\medskip\noindent
Proof of (8). Lemma \href{https://stacks.math.columbia.edu/tag/00TP}{lemma finite presentation fs Noetherian (Tag 00TP)}
gives the result for smooth ring maps. The resulting smooth ring map
$R_0 \to S_0$ satisfies the
hypotheses of Lemma \href{https://stacks.math.columbia.edu/tag/00RL}{lemma relative dimension CM (Tag 00RL)}, and hence we may
replace $S_0$ by the factor of relative dimension $0$ over $R_0$.

\medskip\noindent
Proof of (9). Follows from (8) since $R_0 \to A$ will factor through
$A_i$ for some $i$ by Lemma \href{https://stacks.math.columbia.edu/tag/00QO}{lemma characterize finite presentation (Tag 00QO)}.

\medskip\noindent
Proof of (10). Follows from (9), (1), and (2) since $S^{-1}A$ is a
filtered colimit of principal localizations of $A$.

\medskip\noindent
Proof of (11). Use Lemma \href{https://stacks.math.columbia.edu/tag/0GIF}{lemma product smooth (Tag 0GIF)} to see the
result for smoothness and then use that $\Omega_{B/A}$
is zero if and only if both $\Omega_{B'/A}$ and $\Omega_{B''/A}$
are zero.
\end{proof}


\begin{proposition}
\label{proposition-etale-locally-standard}
Let $R \to S$ be a ring map. Let $\mathfrak q \subset S$ be a prime.
If $R \to S$ is \'etale at $\mathfrak q$, then there exists
a $g \in S$, $g \not \in \mathfrak q$ such that $R \to S_g$
is standard \'etale.
\end{proposition}

\begin{proof}
The following proof is a little roundabout and there may be ways to
shorten it.

\medskip\noindent
Step 1. By Definition \ref{definition-etale}
there exists a $g \in S$, $g \not \in \mathfrak q$
such that $R \to S_g$ is \'etale. Thus we may assume that $S$ is \'etale
over $R$.

\medskip\noindent
Step 2. By Lemma \ref{lemma-etale} there exists an \'etale ring map
$R_0 \to S_0$ with $R_0$ of finite type over $\mathbf{Z}$, and a ring map
$R_0 \to R$ such that $R = R \otimes_{R_0} S_0$. Denote
$\mathfrak q_0$ the prime of $S_0$ corresponding to $\mathfrak q$.
If we show the result for $(R_0 \to S_0, \mathfrak q_0)$ then the
result follows for $(R \to S, \mathfrak q)$ by base change. Hence
we may assume that $R$ is Noetherian.

\medskip\noindent
Step 3.
Note that $R \to S$ is quasi-finite by Lemma \href{https://stacks.math.columbia.edu/tag/00U5}{lemma etale quasi finite (Tag 00U5)}.
By Lemma \href{https://stacks.math.columbia.edu/tag/00QB}{lemma quasi finite open integral closure (Tag 00QB)}
there exists a finite ring map $R \to S'$, an $R$-algebra map
$S' \to S$, an element $g' \in S'$ such that
$g' \not \in \mathfrak q$ such that $S' \to S$ induces
an isomorphism $S'_{g'} \cong S_{g'}$.
(Note that of course $S'$ is not \'etale over $R$ in general.)
Thus we may assume that (a) $R$ is Noetherian, (b) $R \to S$ is finite
and (c) $R \to S$ is \'etale at $\mathfrak q$
(but no longer necessarily \'etale at all primes).

\medskip\noindent
Step 4. Let $\mathfrak p \subset R$ be the prime corresponding
to $\mathfrak q$. Consider the fibre ring
$S \otimes_R \kappa(\mathfrak p)$. This is a finite algebra over
$\kappa(\mathfrak p)$. Hence it is Artinian
(see Lemma \href{https://stacks.math.columbia.edu/tag/00J6}{lemma finite dimensional algebra (Tag 00J6)}) and
so a finite product of local rings
$$
S \otimes_R \kappa(\mathfrak p) = \prod\nolimits_{i = 1}^n A_i
$$
see Proposition \href{https://stacks.math.columbia.edu/tag/00KJ}{proposition dimension zero ring (Tag 00KJ)}. One of the factors,
say $A_1$, is the local ring $S_{\mathfrak q}/\mathfrak pS_{\mathfrak q}$
which is isomorphic to $\kappa(\mathfrak q)$,
see Lemma \href{https://stacks.math.columbia.edu/tag/00U4}{lemma etale at prime (Tag 00U4)}. The other factors correspond to
the other primes, say $\mathfrak q_2, \ldots, \mathfrak q_n$ of
$S$ lying over $\mathfrak p$.

\medskip\noindent
Step 5. We may choose a nonzero element $\alpha \in \kappa(\mathfrak q)$ which
generates the finite separable field extension
$\kappa(\mathfrak q)/\kappa(\mathfrak p)$ (so even if the
field extension is trivial we do not allow $\alpha = 0$).
Note that for any $\lambda \in \kappa(\mathfrak p)^*$ the
element $\lambda \alpha$ also generates $\kappa(\mathfrak q)$
over $\kappa(\mathfrak p)$. Consider the element
$$
\overline{t} =
(\alpha, 0, \ldots, 0) \in
\prod\nolimits_{i = 1}^n A_i =
S \otimes_R \kappa(\mathfrak p).
$$
After possibly replacing $\alpha$ by $\lambda \alpha$ as above
we may assume that $\overline{t}$ is the image of $t \in S$.
Let $I \subset R[x]$ be the kernel of the $R$-algebra
map $R[x] \to S$ which maps $x$ to $t$. Set $S' = R[x]/I$,
so $S' \subset S$. Here is a diagram
$$
\xymatrix{
R[x] \ar[r] & S' \ar[r] & S \\
R \ar[u] \ar[ru] \ar[rru] & &
}
$$
By construction the primes $\mathfrak q_j$, $j \geq 2$ of $S$ all
lie over the prime $(\mathfrak p, x)$ of $R[x]$, whereas
the prime $\mathfrak q$ lies over a different prime of $R[x]$
because $\alpha \not = 0$.

\medskip\noindent
Step 6. Denote $\mathfrak q' \subset S'$ the prime of $S'$
corresponding to $\mathfrak q$. By the above $\mathfrak q$ is
the only prime of $S$ lying over $\mathfrak q'$. Thus we see that
$S_{\mathfrak q} = S_{\mathfrak q'}$, see
Lemma \href{https://stacks.math.columbia.edu/tag/00EA}{lemma unique prime over localize below (Tag 00EA)} (we have
going up for $S' \to S$ by Lemma \href{https://stacks.math.columbia.edu/tag/00GU}{lemma integral going up (Tag 00GU)}
since $S' \to S$ is finite as $R \to S$ is finite).
It follows that $S'_{\mathfrak q'} \to S_{\mathfrak q}$ is finite
and injective as the localization of the finite injective ring map
$S' \to S$. Consider the maps of local rings
$$
R_{\mathfrak p} \to S'_{\mathfrak q'} \to S_{\mathfrak q}
$$
The second map is finite and injective. We have
$S_{\mathfrak q}/\mathfrak pS_{\mathfrak q} = \kappa(\mathfrak q)$,
see Lemma \href{https://stacks.math.columbia.edu/tag/00U4}{lemma etale at prime (Tag 00U4)}.
Hence a fortiori
$S_{\mathfrak q}/\mathfrak q'S_{\mathfrak q} = \kappa(\mathfrak q)$.
Since
$$
\kappa(\mathfrak p) \subset \kappa(\mathfrak q') \subset \kappa(\mathfrak q)
$$
and since $\alpha$ is in the image of $\kappa(\mathfrak q')$ in
$\kappa(\mathfrak q)$
we conclude that $\kappa(\mathfrak q') = \kappa(\mathfrak q)$.
Hence by Nakayama's Lemma \href{https://stacks.math.columbia.edu/tag/00DV}{lemma NAK (Tag 00DV)} applied to the
$S'_{\mathfrak q'}$-module map $S'_{\mathfrak q'} \to S_{\mathfrak q}$,
the map $S'_{\mathfrak q'} \to S_{\mathfrak q}$ is surjective.
In other words,
$S'_{\mathfrak q'} \cong S_{\mathfrak q}$.

\medskip\noindent
Step 7. By Lemma \href{https://stacks.math.columbia.edu/tag/00QS}{lemma isomorphic local rings (Tag 00QS)} there exist
$g \in S$, $g \not \in \mathfrak q$ and $g' \in S'$, $g' \not \in \mathfrak q'$
such that $S'_{g'} \cong S_g$. As $R$ is Noetherian the ring $S'$ is finite
over $R$ because it is an $R$-submodule
of the finite $R$-module $S$. Hence after replacing $S$ by $S'$ we may
assume that (a) $R$ is Noetherian, (b) $S$ finite over $R$, (c)
$S$ is \'etale over $R$ at $\mathfrak q$, and (d) $S = R[x]/I$.

\medskip\noindent
Step 8. Consider the ring
$S \otimes_R \kappa(\mathfrak p) = \kappa(\mathfrak p)[x]/\overline{I}$
where $\overline{I} = I \cdot \kappa(\mathfrak p)[x]$ is the ideal generated
by $I$ in $\kappa(\mathfrak p)[x]$. As $\kappa(\mathfrak p)[x]$ is a PID
we know that $\overline{I} = (\overline{h})$ for some monic
$\overline{h} \in \kappa(\mathfrak p)[x]$. After replacing $\overline{h}$
by $\lambda \cdot \overline{h}$ for some $\lambda \in \kappa(\mathfrak p)$
we may assume that $\overline{h}$ is the image of some $h \in I \subset R[x]$.
(The problem is that we do not know if we may choose $h$ monic.)
Also, as in Step 4 we know that
$S \otimes_R \kappa(\mathfrak p) = A_1 \times \ldots \times A_n$ with
$A_1 = \kappa(\mathfrak q)$ a finite separable extension of
$\kappa(\mathfrak p)$ and $A_2, \ldots, A_n$ local. This implies
that
$$
\overline{h} = \overline{h}_1 \overline{h}_2^{e_2} \ldots \overline{h}_n^{e_n}
$$
for certain pairwise coprime irreducible monic polynomials
$\overline{h}_i \in \kappa(\mathfrak p)[x]$ and certain
$e_2, \ldots, e_n \geq 1$. Here the numbering is chosen so that
$A_i = \kappa(\mathfrak p)[x]/(\overline{h}_i^{e_i})$ as
$\kappa(\mathfrak p)[x]$-algebras. Note that $\overline{h}_1$ is
the minimal polynomial of $\alpha \in \kappa(\mathfrak q)$ and hence
is a separable polynomial (its derivative is prime to itself).

\medskip\noindent
Step 9. Let $m \in I$ be a monic element; such an element exists
because the ring extension $R \to R[x]/I$ is finite hence integral.
Denote $\overline{m}$ the image in $\kappa(\mathfrak p)[x]$.
We may factor
$$
\overline{m} = \overline{k}
\overline{h}_1^{d_1} \overline{h}_2^{d_2} \ldots \overline{h}_n^{d_n}
$$
for some $d_1 \geq 1$, $d_j \geq e_j$, $j = 2, \ldots, n$ and
$\overline{k} \in \kappa(\mathfrak p)[x]$ prime to all the $\overline{h}_i$.
Set $f = m^l + h$ where $l \deg(m) > \deg(h)$, and $l \geq 2$.
Then $f$ is monic as a polynomial over $R$. Also, the image $\overline{f}$
of $f$ in $\kappa(\mathfrak p)[x]$ factors as
$$
\overline{f} =
\overline{h}_1 \overline{h}_2^{e_2} \ldots \overline{h}_n^{e_n}
+
\overline{k}^l \overline{h}_1^{ld_1} \overline{h}_2^{ld_2}
\ldots \overline{h}_n^{ld_n}
=
\overline{h}_1(\overline{h}_2^{e_2} \ldots \overline{h}_n^{e_n}
+
\overline{k}^l
\overline{h}_1^{ld_1 - 1} \overline{h}_2^{ld_2} \ldots \overline{h}_n^{ld_n})
= \overline{h}_1 \overline{w}
$$
with $\overline{w}$ a polynomial relatively prime to $\overline{h}_1$.
Set $g = f'$ (the derivative with respect to $x$).

\medskip\noindent
Step 10. The ring map $R[x] \to S = R[x]/I$ has the properties:
(1) it maps $f$ to zero, and
(2) it maps $g$ to an element of $S \setminus \mathfrak q$.
The first assertion is clear since $f$ is an element of $I$.
For the second assertion we just have to show that $g$ does
not map to zero in
$\kappa(\mathfrak q) = \kappa(\mathfrak p)[x]/(\overline{h}_1)$.
The image of $g$ in $\kappa(\mathfrak p)[x]$ is the derivative
of $\overline{f}$. Thus (2) is clear because
$$
\overline{g} =
\frac{\text{d}\overline{f}}{\text{d}x} =
\overline{w}\frac{\text{d}\overline{h}_1}{\text{d}x} +
\overline{h}_1\frac{\text{d}\overline{w}}{\text{d}x},
$$
$\overline{w}$ is prime to $\overline{h}_1$ and
$\overline{h}_1$ is separable.

\medskip\noindent
Step 11.
We conclude that $\varphi : R[x]/(f) \to S$ is a surjective ring map,
$R[x]_g/(f)$ is \'etale over $R$ (because it is standard \'etale,
see Lemma \href{https://stacks.math.columbia.edu/tag/00UC}{lemma standard etale (Tag 00UC)}) and $\varphi(g) \not \in \mathfrak q$.
Pick an element $g' \in R[x]/(f)$ such that
also $\varphi(g') \not \in \mathfrak q$ and $S_{\varphi(g')}$
is \'etale over $R$ (which exists since $S$ is \'etale over $R$ at
$\mathfrak q$). Then the ring map
$R[x]_{gg'}/(f) \to S_{\varphi(gg')}$ is a surjective map of \'etale
algebras over $R$. Hence it is \'etale by Lemma \href{https://stacks.math.columbia.edu/tag/00U7}{lemma map between etale (Tag 00U7)}.
Hence it is a localization by
Lemma \href{https://stacks.math.columbia.edu/tag/00U8}{lemma surjective flat finitely presented (Tag 00U8)}.
Thus a localization of $S$ at an element not in $\mathfrak q$ is
isomorphic to a localization of a standard \'etale algebra over $R$
which is what we wanted to show.
\end{proof}
\end{document}
